\section*{Ethics Statement}

This work studies whether and how well LLMs can improve the code of other LLM
agents. Automating agent improvement carries a dual-use tension, since the same
capability that repairs and strengthens a benign agent could in principle be
turned toward a harmful one. The benchmark's design limits this exposure. The
optimization target is always a bounded, benign downstream task (document QA,
deep research, multi-step reasoning, or terminal use), the score is the task's
own verifier on a held-out split and nothing else, and every optimizer runs inside
a trusted harness that meters spend, versions every change for audit, and
confines the target to a fixed model behind an allow-list -- a concrete
precaution given evidence that comparable agent benchmarks can be gamed via
evaluator hijacking or judge prompt-injection \citep{benchjack2026}. Nothing in
the setup rewards or requires acquiring new capabilities, tools, or resources
beyond editing a fixed agent's code.

Our benchmark reuses existing public datasets and benchmarks under their
respective licenses and adds seed agents and split definitions that we release.
The downstream corpora are public documents, and we introduce no personal or
sensitive data. Because agent evaluation is compute-intensive, we report token
usage as a first-class metric to make the cost of this line of work visible, and
we keep the search budget denominated in evaluation calls and case-runs so that comparisons do not
implicitly favor optimizers with larger compute budgets. Finally, we report
the completed evaluation, including repeated runs, with explicit caveats (see
the Limitations section) and avoid over-claiming fine-grained model rankings.
